\documentclass[12pt]{article}
\usepackage{fullpage}
\usepackage[T1]{fontenc}
\usepackage[utf8]{inputenc}
\usepackage{lmodern}
\usepackage{microtype}
\usepackage{amsmath,amssymb,amsthm}
\usepackage{hyperref}

\newtheorem{theorem}{Theorem}
\newtheorem{lemma}{Lemma}
\newtheorem{proposition}{Proposition}
\newtheorem{corollary}{Corollary}
\theoremstyle{definition}
\newtheorem{remark}{Remark}

\newcommand{\R}{\mathbb{R}}

\title{Global stability of the positive equilibrium for autonomous and
nonautonomous Nicholson blowflies equations}
\author{}
\date{}

\begin{document}
\maketitle

\section{Setup, standing assumptions, and notation}

Throughout, $P>d>0$ and $s>0$ are fixed. Write
\[
  f(x)=x e^{-x}\qquad(x\ge 0),\qquad
  N^{*}=\ln\frac{P}{d}\;(>0),\qquad r:=\frac{P}{d}=e^{N^{*}}>1 .
\]
The autonomous equation is
\begin{equation}\label{eq:auto}
  N'(t)=-d\,N(t)+P\,f\bigl(N(t-s)\bigr),\qquad t\ge 0 .
\end{equation}
The three nonautonomous models on $t\ge0$ are
\begin{align}
  N'(t)&=-d\,N(t)+P\!\int_{h(t)}^{t}\! f\bigl(N(u)\bigr)\,d_uR(t,u),\label{eq:A}\tag{A}\\
  N'(t)&=-d\,N(t)+P\sum_{k=1}^{m}a_k(t)\,f\bigl(N(h_k(t))\bigr),\label{eq:B}\tag{B}\\
  N'(t)&=-d\,N(t)+P\!\int_{h(t)}^{t}\! K(t,u)\,f\bigl(N(u)\bigr)\,du,\label{eq:C}\tag{C}
\end{align}
under the standing assumptions of the problem statement:
\begin{itemize}
\item $h,h_k$ are measurable, $t-s\le h(t)\le t$ and $t-s\le h_k(t)\le t$;
\item for each $t$, $R(t,\cdot)$ is nondecreasing of bounded variation with
$R(t,t)-R(t,h(t))=1$;
\item $a_k(t)\ge 0$ and $\sum_{k=1}^m a_k(t)\equiv 1$;
\item $K(t,u)\ge 0$ and $\int_{h(t)}^t K(t,u)\,du=1$.
\end{itemize}
We consider solutions with continuous initial data
$N|_{[-s,0]}=\varphi\in C([-s,0],\R)$, $\varphi\ge 0$, $\varphi(0)>0$.
Such solutions exist, are unique on $[0,\infty)$, and remain positive: since
$f\ge0$ we have $N'(t)\ge -dN(t)$, so $N(t)\ge N(0)e^{-dt}>0$ for all $t\ge0$.
A direct substitution shows $N^{*}=\ln(P/d)$ is an equilibrium of
\eqref{eq:auto}--\eqref{eq:C}, since $Pf(N^*)=PN^*e^{-N^*}=dN^*$ (using
$Pe^{-N^*}=d$).

\paragraph{Unified form.}
We write every model as
\begin{equation}\label{eq:unified}
  N'(t)=-d\,N(t)+P\,\bigl(\mathcal A_t f(N)\bigr),
\end{equation}
where, for a bounded function $\psi$ on $[t-s,t]$,
\[
  \mathcal A_t\psi=\int_{h(t)}^{t}\psi(u)\,d_uR(t,u),\quad
  \sum_{k=1}^m a_k(t)\psi(h_k(t)),\quad
  \int_{h(t)}^{t}K(t,u)\psi(u)\,du
\]
for \eqref{eq:A},\eqref{eq:B},\eqref{eq:C} respectively, and
$\mathcal A_t\psi=\psi(t-s)$ for \eqref{eq:auto}. In all cases $\mathcal A_t$ is
a \emph{positive averaging operator over the window $[t-s,t]$}:
\begin{itemize}
\item[(P1)] (monotone) if $\psi_1\le\psi_2$ on $[t-s,t]$ then
  $\mathcal A_t\psi_1\le\mathcal A_t\psi_2$;
\item[(P2)] (reproduces constants) $\mathcal A_t c=c$, since the total weight
  mass is $1$;
\item[(P3)] (enclosure) consequently
  $\displaystyle\inf_{u\in[t-s,t]}\psi(u)\le\mathcal A_t\psi
  \le\sup_{u\in[t-s,t]}\psi(u)$.
\end{itemize}
\emph{Every estimate below uses only (P1)--(P3).} Hence each result holds
simultaneously for \eqref{eq:auto} and for \eqref{eq:A}--\eqref{eq:C}, with no
restriction on the delay beyond $t-s\le h(\cdot)\le t$.

$N^*$ is a \emph{global attractor} (within the positive cone) if
$N(t)\to N^*$ for every solution with $\varphi\ge0,\varphi(0)>0$; it is
\emph{globally asymptotically stable} if in addition it is Lyapunov stable.

We use the elementary properties of $f(x)=xe^{-x}$: $f(0)=0$; $f$ strictly
increasing on $[0,1]$ and strictly decreasing on $[1,\infty)$;
$\max_{x\ge0}f(x)=f(1)=e^{-1}$; $f(x)>0$ for $x>0$.

\section{Dissipativity (permanence)}

\begin{lemma}[Uniform persistence and boundedness]\label{lem:perm}
Let $N$ be any solution of \eqref{eq:unified} with $\varphi\ge0$,
$\varphi(0)>0$. Set $M_0:=Pe^{-1}/d=re^{-1}$ and $B:=M_0+1$, and fix any
\begin{equation}\label{eq:etachoice}
  \eta\in\Bigl(0,\ \min\{1,\ N^*,\ Pf(B)/d\}\Bigr).
\end{equation}
Then
\begin{equation}\label{eq:limsup}
  \limsup_{t\to\infty}N(t)\le M_0,
\end{equation}
and there is $T_*$ (depending on the solution) with
\begin{equation}\label{eq:trap}
  \eta\le N(t)\le B\qquad(t\ge T_*).
\end{equation}
In particular $\liminf_{t\to\infty}N(t)\ge\eta$, where $\eta$ depends only on
$P,d$.
\end{lemma}

\begin{proof}
\emph{Upper bound.} By (P3) and $f\le e^{-1}$, $\mathcal A_tf(N)\le e^{-1}$,
hence $N'(t)\le -dN(t)+Pe^{-1}$. Comparison with $y'=-dy+Pe^{-1}$, whose
solutions converge to $Pe^{-1}/d=M_0$, gives a time $T_1$ with $N(t)\le B=M_0+1$
for $t\ge T_1$, and the sharper bound \eqref{eq:limsup} (the comparison solution
tends to $M_0$, so $\limsup N\le M_0$).

\emph{Lower bound (persistence via a fluctuation argument).} By the choice
\eqref{eq:etachoice},
\begin{equation}\label{eq:cetabound}
  Pf(\eta)=P\eta e^{-\eta}=d\eta\,re^{-\eta}>d\eta
  \quad(\text{since }\eta<N^*=\ln r),\qquad
  Pf(B)>d\eta\quad(\text{since }\eta<Pf(B)/d).
\end{equation}

We show $\ell:=\liminf_{t\to\infty}N(t)\ge\eta$. For $t\ge T_1$ the
derivative is bounded, $|N'(t)|\le dB+Pe^{-1}$ (as $N\le B$, $f\le e^{-1}$),
so $N$ is uniformly continuous on $[T_1,\infty)$. By the fluctuation lemma there
is a sequence $t_n\to\infty$, $t_n\ge T_1+s$, with
$N(t_n)\to\ell$, $N(t_n)=\min_{[T_1,t_n]}N$, and $N'(t_n)\le0$. At such a
record-minimum time, $N(u)\ge N(t_n)$ for all $u\in[t_n-s,t_n]$, and
$N(u)\le B$; hence
$f(N(u))\ge\min\{f(N(t_n)),f(B)\}$ for each such $u$ (use $f$ increasing on
$[0,1]$ when $N(u)\le1$, and $f$ decreasing on $[1,\infty)$ when $N(u)>1$). By
(P1)--(P2), $\mathcal A_{t_n}f(N)\ge\min\{f(N(t_n)),f(B)\}$, so from
$N'(t_n)\le0$,
\begin{equation}\label{eq:recordmin}
  d\,N(t_n)\ge P\,\mathcal A_{t_n}f(N)\ge P\min\{f(N(t_n)),\,f(B)\}.
\end{equation}
Suppose $\ell<\eta$. Since $N(t_n)\to\ell<\eta\le1$, for all large $n$ we have
$N(t_n)<\eta$, hence $f(N(t_n))=N(t_n)e^{-N(t_n)}$ and, because
$Pf(B)>d\eta>d\,N(t_n)$ by \eqref{eq:cetabound}, the minimum in
\eqref{eq:recordmin} is $f(N(t_n))$ (otherwise \eqref{eq:recordmin} would read
$d\,N(t_n)\ge Pf(B)>d\,N(t_n)$, impossible). Thus for large $n$,
\[
  d\,N(t_n)\ge P\,N(t_n)e^{-N(t_n)} .
\]
Dividing by $N(t_n)>0$ gives $d\ge Pe^{-N(t_n)}=d\,r\,e^{-N(t_n)}$, i.e.
$re^{-N(t_n)}\le1$, i.e. $N(t_n)\ge\ln r=N^*$. But $N(t_n)\to\ell<\eta<N^*$,
a contradiction for large $n$. Therefore $\ell\ge\eta$. Combined with
$\limsup N\le M_0\le B$, there is $T_*$ with $N(t)\in[\eta,B]$ for $t\ge T_*$,
proving \eqref{eq:trap}.

\end{proof}

\section{The one-dimensional reduction and the contraction inequality}

The asymptotics of \eqref{eq:unified} are governed by
\begin{equation}\label{eq:G}
  G(x):=\frac{P}{d}\,f(x)=r\,x\,e^{-x},\qquad x\ge0 ,
\end{equation}
with unique positive fixed point $G(N^*)=N^*$ (as $G(x)=x\iff re^{-x}=1\iff
x=N^*$), $G(0)=0$, $G$ increasing on $[0,1]$, decreasing on $[1,\infty)$,
$\max G=G(1)=r/e=M_0$, and
\begin{equation}\label{eq:Gprime}
  G'(x)=r(1-x)e^{-x},\qquad G'(N^*)=1-N^* .
\end{equation}

\begin{lemma}[Sharp contraction toward $N^*$]\label{lem:contr}
Let $N^*=\ln r$. Then
\begin{equation}\label{eq:contr}
  |G(x)-N^*|<|x-N^*|\qquad\text{for all }x>0,\ x\ne N^*,
\end{equation}
\emph{if and only if} $N^*<2$. Moreover, when $N^*<2$, on every compact
interval $[m,M]\subset(0,\infty)$ the contraction is uniform:
\[
  q:=\sup_{x\in[m,M],\,x\ne N^*}\frac{|G(x)-N^*|}{|x-N^*|}<1 .
\]
\end{lemma}

\begin{proof}
\eqref{eq:contr} is equivalent to
$\bigl(G(x)-N^*\bigr)^2-\bigl(x-N^*\bigr)^2<0$, i.e.
\[
  \bigl(G(x)-x\bigr)\bigl(G(x)+x-2N^*\bigr)<0 .
\]
Define $A(x):=G(x)-x=x(e^{N^*-x}-1)$ and $B(x):=G(x)+x-2N^*$. For $x>0$,
$\operatorname{sgn}A(x)=\operatorname{sgn}(N^*-x)$ (with $A(N^*)=0$), so
\eqref{eq:contr} holds at $x\ne N^*$ iff $A(x)B(x)<0$, i.e. iff
$\operatorname{sgn}B(x)=\operatorname{sgn}(x-N^*)$.

Now $B(N^*)=N^*+N^*-2N^*=0$, and
\[
  B'(x)=G'(x)+1=r(1-x)e^{-x}+1=(1-x)e^{N^*-x}+1 .
\]
The map $x\mapsto(1-x)e^{-x}$ has derivative $(x-2)e^{-x}$, vanishing only at
$x=2$, where it attains its global minimum $(1-2)e^{-2}=-e^{-2}$; hence
$(1-x)e^{-x}\ge -e^{-2}$ for all $x$, giving
\begin{equation}\label{eq:Bplow}
  B'(x)=e^{N^*}\,(1-x)e^{-x}+1\ \ge\ 1-e^{N^*-2}\qquad(x>0).
\end{equation}

\emph{If $N^*<2$:} $e^{N^*-2}<1$, so $B'(x)\ge 1-e^{N^*-2}>0$ for all $x>0$;
thus $B$ is strictly increasing and, with $B(N^*)=0$,
$\operatorname{sgn}B(x)=\operatorname{sgn}(x-N^*)$ for $x\ne N^*$. Hence
$A(x)B(x)<0$, proving \eqref{eq:contr}.

\emph{If $N^*\ge2$:} then $B'(N^*)=(1-N^*)+1=2-N^*\le0$. Since $B(N^*)=0$, for
$x$ slightly above $N^*$ we have $B(x)\le0$ while $A(x)<0$, so $A(x)B(x)\ge0$
and $|G(x)-N^*|\ge|x-N^*|$ there. Thus \eqref{eq:contr} fails; this proves the
``iff''.

\emph{Uniformity.} Assume $N^*<2$. Set $\rho(x):=|G(x)-N^*|/|x-N^*|$ for
$x\ne N^*$, extended by $\rho(N^*):=|G'(N^*)|=|1-N^*|<1$ (the limit, by
\eqref{eq:Gprime} and l'Hôpital). Then $\rho$ is continuous on $[m,M]$, with
$\rho<1$ everywhere by \eqref{eq:contr} (and $\rho(N^*)<1$). A continuous
function that is $<1$ on a compact set attains a maximum $q<1$.
\end{proof}

\begin{lemma}[Collapse of the invariant interval iteration]\label{lem:iter}
Assume $0<N^*<2$. Fix $0<m_0<N^*$ and $M_0\ge N^*$, and define
\begin{equation}\label{eq:Mm}
  M_{k+1}=\max_{x\in[m_k,M_k]}G(x),\qquad
  m_{k+1}=\min_{x\in[m_k,M_k]}G(x)\qquad(k\ge0).
\end{equation}
Then $N^*\in[m_{k+1},M_{k+1}]\subseteq[m_k,M_k]$ for all $k$, and
$M_k\downarrow N^*$, $m_k\uparrow N^*$.
\end{lemma}

\begin{proof}
Let $q<1$ be the uniform contraction constant of Lemma~\ref{lem:contr} on the
fixed compact interval $[m_0,M_0]$.

We prove by induction that $N^*\in[m_k,M_k]$ and
$[m_{k+1},M_{k+1}]\subseteq[m_k,M_k]$, with
\begin{equation}\label{eq:Dk}
  D_k:=\max\{M_k-N^*,\ N^*-m_k\}\ \text{satisfies}\ D_{k+1}\le qD_k .
\end{equation}
Base $k=0$: $m_0<N^*\le M_0$ by hypothesis, so $N^*\in[m_0,M_0]$.
Inductive step: assume $N^*\in[m_k,M_k]\subseteq[m_0,M_0]$. Since $G(N^*)=N^*$
and $N^*\in[m_k,M_k]$, the range $G([m_k,M_k])$ contains $N^*$, whence
$m_{k+1}\le N^*\le M_{k+1}$. For any $x\in[m_k,M_k]$ we have $|x-N^*|\le D_k$,
so by Lemma~\ref{lem:contr} (uniform constant $q$ on $[m_0,M_0]$)
$|G(x)-N^*|\le q\,|x-N^*|\le qD_k$; therefore
$G([m_k,M_k])\subseteq[N^*-qD_k,\,N^*+qD_k]$, i.e.
\[
  M_{k+1}-N^*\le qD_k,\qquad N^*-m_{k+1}\le qD_k,
\]
which gives $D_{k+1}=\max\{M_{k+1}-N^*,\,N^*-m_{k+1}\}\le qD_k$, proving
\eqref{eq:Dk}. Since $q<1$, $D_k\le qD_{k-1}$ shows in particular
$D_{k+1}\le D_k$, so $M_{k+1}-N^*\le M_k-N^*$ and $N^*-m_{k+1}\le N^*-m_k$;
equivalently $m_k\le m_{k+1}\le N^*\le M_{k+1}\le M_k$, i.e.
$[m_{k+1},M_{k+1}]\subseteq[m_k,M_k]$ and the sequences are monotone (with
$m_k\uparrow$, $M_k\downarrow$). Iterating \eqref{eq:Dk} gives
$D_k\le q^kD_0\to0$, hence $M_k\to N^*$ and $m_k\to N^*$.
\end{proof}

\section{(i)--(ii) Sufficient global stability, uniform in the delay structure}

\begin{theorem}[Delay-independent global stability]\label{thm:main}
Suppose $P>d>0$ and
\begin{equation}\label{eq:cond}
  N^*=\ln\frac{P}{d}<2,\qquad\text{equivalently}\qquad d<P<de^{2}.
\end{equation}
Then for every admissible delay structure (any measurable $h,h_k$ with values
in $[t-s,t]$, any $R,a_k,K$ as specified, any $s>0$), the positive equilibrium
$N^*$ of each of \eqref{eq:auto}, \eqref{eq:A}, \eqref{eq:B}, \eqref{eq:C} is a
global attractor: $N(t)\to N^*$ for every solution with $\varphi\ge0$,
$\varphi(0)>0$.
\end{theorem}

\begin{proof}
Use the unified form \eqref{eq:unified}; every step uses only (P1)--(P3).

Set $\widehat M:=re^{-1}+1$ and $B:=\widehat M$. By Lemma~\ref{lem:perm},
fix $\eta\in\bigl(0,\min\{1,N^*,Pf(B)/d\}\bigr)$ and let $m_0:=\eta$,
$M_0:=\widehat M$. By \eqref{eq:trap} there is $T_0$ with
$N(t)\in[m_0,M_0]$ for all $t\ge T_0$. Note $m_0<N^*$ and $M_0=re^{-1}+1>re^{-1}
=\max G\ge G(N^*)=N^*$, so the pair $(m_0,M_0)$ satisfies the hypotheses
$0<m_0<N^*\le M_0$ of Lemma~\ref{lem:iter}. Define
\[
  L:=\limsup_{t\to\infty}N(t),\qquad \ell:=\liminf_{t\to\infty}N(t),
  \qquad m_0\le\ell\le L\le M_0 .
\]

\emph{Step 1 (one comparison round).} Fix $\varepsilon>0$; there is $T_1\ge T_0$
with
\begin{equation}\label{eq:enclose}
  \ell-\varepsilon\le N(u)\le L+\varepsilon\qquad(u\ge T_1).
\end{equation}
For $t\ge T_1+s$ all delayed arguments $u\in[t-s,t]$ obey \eqref{eq:enclose}, so
$f(N(u))\in\bigl[\min_{[\ell-\varepsilon,L+\varepsilon]}f,
\max_{[\ell-\varepsilon,L+\varepsilon]}f\bigr]$, and by (P1),(P3),
\[
  \min_{[\ell-\varepsilon,L+\varepsilon]}f\le\mathcal A_tf(N)
  \le\max_{[\ell-\varepsilon,L+\varepsilon]}f .
\]
With $G=(P/d)f$, $\overline M_\varepsilon:=\max_{[\ell-\varepsilon,L+\varepsilon]}G$,
$\underline m_\varepsilon:=\min_{[\ell-\varepsilon,L+\varepsilon]}G$, this gives
\begin{equation}\label{eq:diffineq}
  -d\bigl(N(t)-\underline m_\varepsilon\bigr)\le N'(t)\le
  -d\bigl(N(t)-\overline M_\varepsilon\bigr),\qquad t\ge T_1+s .
\end{equation}

\emph{Step 2 ($\limsup/\liminf$).} The right inequality is
$N'\le -d(N-\overline M_\varepsilon)$; comparison with
$y'=-d(y-\overline M_\varepsilon)$ (solutions $\to\overline M_\varepsilon$)
yields $L\le\overline M_\varepsilon$. The left inequality similarly gives
$\ell\ge\underline m_\varepsilon$. Letting $\varepsilon\downarrow0$ and using
continuity of $G$,
\begin{equation}\label{eq:LL}
  L\le \max_{x\in[\ell,L]}G(x),\qquad \ell\ge\min_{x\in[\ell,L]}G(x).
\end{equation}

\emph{Step 3 (trapping in the iteration).} Let $(m_k,M_k)$ be the iteration of
Lemma~\ref{lem:iter} with the same $m_0,M_0$. By induction
$[\ell,L]\subseteq[m_k,M_k]$: the base $k=0$ is $m_0\le\ell\le L\le M_0$; if
$[\ell,L]\subseteq[m_k,M_k]$, then $\max_{[\ell,L]}G\le\max_{[m_k,M_k]}G=M_{k+1}$
and $\min_{[\ell,L]}G\ge\min_{[m_k,M_k]}G=m_{k+1}$, so by \eqref{eq:LL}
$L\le M_{k+1}$ and $\ell\ge m_{k+1}$, i.e.
$[\ell,L]\subseteq[m_{k+1},M_{k+1}]$.

By Lemma~\ref{lem:iter} ($N^*<2$), $[m_k,M_k]\downarrow\{N^*\}$, hence
$[\ell,L]\subseteq\bigcap_k[m_k,M_k]=\{N^*\}$, so $\ell=L=N^*$ and
$N(t)\to N^*$. The argument is identical for all four models, since
$\mathcal A_t$ entered only through (P1)--(P3) over the window $[t-s,t]$.
\end{proof}

\begin{remark}[Strength; answer to (i)]
The threshold $N^*<2$ (i.e. $P/d<e^2$) is exactly the delay-independent local
asymptotic stability boundary of the autonomous equation
(Section~\ref{sec:nec}); thus Theorem~\ref{thm:main} gives, for each of
\eqref{eq:A}--\eqref{eq:C}, a global stability criterion as strong as the best
delay-independent criterion for the autonomous problem, and the constant $2$ is
optimal (at $N^*=2$ local stability fails for large delay). The classical
monotone case $N^*\le1$ ($P\le de$, where eventually $N\in(0,1]$ and $f$ is
increasing) is included.
\end{remark}

\begin{corollary}[Answer to (ii)]\label{cor:ii}
The delay-independent criterion $N^*<2$ of the autonomous equation remains valid
\emph{verbatim} when the constant delay $s$ is replaced by an arbitrary
measurable variable delay $h(t)\in[t-s,t]$, and more generally by any of
\eqref{eq:A}--\eqref{eq:C}: Theorem~\ref{thm:main} imposes no condition on $h$
beyond admissibility. By contrast, the stronger delay-dependent (``$3/2$''-type)
criteria the autonomous equation enjoys for $N^*\ge2$ when $s$ is small do
\emph{not} survive arbitrary variable delay: such criteria require a smallness
bound $ds\le c$, which a large variable delay violates, and—by
Proposition~\ref{prop:nec-auto}—instability genuinely occurs for $N^*>2$ once
the delay exceeds the critical $s_c$. Thus the delay-independent part of the
autonomous criterion is preserved, the delay-smallness part is not.
\end{corollary}

\section{(iii) Necessary delay-dependent conditions}\label{sec:nec}

\subsection{Linearization and characteristic equation}

Let $y(t)=N(t)-N^*$. Using $f'(N^*)=(1-N^*)e^{-N^*}$ and $Pe^{-N^*}=d$, the
linearization of \eqref{eq:auto} about $N^*$ is
\begin{equation}\label{eq:lin}
  y'(t)=-d\,y(t)+d(1-N^*)\,y(t-s),
\end{equation}
and of \eqref{eq:A}--\eqref{eq:C}, $y'(t)=-d\,y(t)+d(1-N^*)\,(\mathcal A_t y)$.

\begin{proposition}[Necessary condition, autonomous]\label{prop:nec-auto}
Set $a:=d>0$, $b:=d(1-N^*)$. The characteristic equation of \eqref{eq:lin} is
\begin{equation}\label{eq:char}
  \lambda+a-b\,e^{-\lambda s}=0 .
\end{equation}
\begin{enumerate}
\item[\rm(a)] If $|b|<a$ (i.e. $0<N^*<2$), all roots have
$\operatorname{Re}\lambda<0$ for every $s\ge0$; the zero solution of
\eqref{eq:lin} is asymptotically stable for all $s$.
\item[\rm(b)] If $N^*>2$ (so $b<-a<0$, $|b|>a$), \eqref{eq:char} has imaginary
roots $\pm i\omega_c$ at
\begin{equation}\label{eq:crit}
  s_c=\frac{\arccos\!\bigl(1/(1-N^*)\bigr)}{d\sqrt{(N^*-1)^2-1}},\qquad
  \omega_c=d\sqrt{(N^*-1)^2-1},
\end{equation}
and a root with $\operatorname{Re}\lambda>0$ for every $s>s_c$; the equilibrium
is unstable for $s>s_c$ (and as $s\to\infty$).
\end{enumerate}
Hence $N^*\le2$ ($P/d\le e^2$) is \emph{necessary} for $N^*$ to be locally
(hence globally) asymptotically stable uniformly in $s$.
\end{proposition}

\begin{proof}
(a) If $\lambda=\sigma+i\omega$, $\sigma\ge0$, solves \eqref{eq:char}, then
$|\lambda+a|=|b|e^{-\sigma s}\le|b|$, while
$|\lambda+a|=\sqrt{(\sigma+a)^2+\omega^2}\ge a>|b|$, a contradiction. Thus all
roots have $\operatorname{Re}\lambda<0$; the strict inequality at $\sigma=0$
prevents accumulation on the imaginary axis, giving asymptotic stability. (At
$s=0$, $\lambda=b-a<0$.)

(b) Put $\lambda=i\omega$, $\omega>0$ in \eqref{eq:char}:
$a=b\cos\omega s$, $\omega=-b\sin\omega s$; squaring and adding gives
$\omega^2+a^2=b^2$, with $\omega_c=\sqrt{b^2-a^2}>0$ since $|b|>a$. With $b<0$,
$\cos\omega_c s=a/b<0$ and $\sin\omega_c s=-\omega_c/b>0$, solved by
$\omega_c s_c=\arccos(a/b)\in(\tfrac\pi2,\pi)$, yielding \eqref{eq:crit} (here
$a/b=1/(1-N^*)$, $b^2-a^2=d^2[(N^*-1)^2-1]>0$). Transversality: from
$F(\lambda,s)=\lambda+a-be^{-\lambda s}$,
$\frac{d\lambda}{ds}=-F_s/F_\lambda
=-b\lambda e^{-\lambda s}/(1+bse^{-\lambda s})$, and using
$be^{-\lambda s}=\lambda+a=a+i\omega_c$ at the crossing,
$\bigl(\tfrac{d\lambda}{ds}\bigr)^{-1}
=-\tfrac{1}{\lambda(a+i\omega_c)}-\tfrac{s_c}{\lambda}$, whose real part at
$\lambda=i\omega_c$ equals $1/(a^2+\omega_c^2)>0$. Hence
$\frac{d(\operatorname{Re}\lambda)}{ds}\big|_{s_c}>0$: the pair crosses into the
right half-plane, so a root with $\operatorname{Re}\lambda>0$ exists for
$s>s_c$, giving instability.
\end{proof}

\begin{proposition}[Necessary conditions for (A)--(C)]\label{prop:nec-nonauto}
Let $N^*=\ln(P/d)$. For each of \eqref{eq:A}--\eqref{eq:C}:
\begin{enumerate}
\item[\rm(1)] If global asymptotic stability holds for \emph{every} admissible
structure over windows $[t-s,t]$ and every $s>0$, then $N^*\le2$
($P\le de^2$). Each model contains \eqref{eq:auto} with delay $s$:
take $R(t,u)=\mathbf1_{\{u\ge t-s\}}$ in \eqref{eq:A}; $m=1,a_1\equiv1,
h_1(t)=t-s$ in \eqref{eq:B}; $K(t,\cdot)\to\delta_{t-s}$ (weak-$*$) in
\eqref{eq:C}. By Proposition~\ref{prop:nec-auto}(b), $N^*>2$ forces instability
for $s>s_c$.
\item[\rm(2)] If the structure realizes a point delay $\tau\in(0,s]$
($h\equiv t-\tau$, $h_k=t-\tau$, or $R/K$ concentrated at $t-\tau$) and
$N^*>2$, then $N^*$ is unstable whenever $\tau>s_c$. Thus a necessary condition
for stability is the delay-dependent bound
\begin{equation}\label{eq:necbound}
  N^*\le2,\quad\text{or}\quad
  N^*>2\ \text{and}\ \tau<s_c
  =\frac{\arccos\!\bigl(1/(1-N^*)\bigr)}{d\sqrt{(N^*-1)^2-1}}\,,
\end{equation}
i.e. $d\tau<\arccos\!\bigl(1/(1-N^*)\bigr)/\sqrt{(N^*-1)^2-1}$ when $N^*>2$.
\end{enumerate}
\end{proposition}

\begin{proof}
(1) Each choice reduces the model to \eqref{eq:auto} with delay $s$ (for
\eqref{eq:C} via $K_n\to\delta_{t-s}$ weak-$*$; the rightmost characteristic
root depends continuously on the kernel in the weak-$*$ topology, so an unstable
root for the point delay persists for $K_n$, $n$ large). By
Proposition~\ref{prop:nec-auto}(b), $N^*>2$ gives a positive-real-part root for
$s>s_c$, contradicting stability; hence $N^*\le2$.

(2) For a structure realizing point delay $\tau$, the linearization is
\eqref{eq:lin} with $s$ replaced by $\tau$; Proposition~\ref{prop:nec-auto}(b)
with delay $\tau$ shows instability when $N^*>2$ and $\tau>s_c$. Stability
requires $\tau\le s_c$, i.e. \eqref{eq:necbound}.
\end{proof}

\begin{remark}[Sharpness]
The delay-independent sufficient condition ($N^*<2$,
Theorem~\ref{thm:main}) and the delay-independent necessary condition
($N^*\le2$, Propositions~\ref{prop:nec-auto}--\ref{prop:nec-nonauto}) coincide
at the boundary value. Hence $N^*=2$ ($P/d=e^2$) is the sharp threshold for
global stability uniform in the delay, for the autonomous equation and for each
of \eqref{eq:A}--\eqref{eq:C}. For $N^*>2$ stability is genuinely
delay-dependent, governed by the explicit critical delay $s_c$ of
\eqref{eq:crit}.
\end{remark}

\section{Summary of answers}

\begin{itemize}
\item[\textbf{(i)}] For each of \eqref{eq:A}, \eqref{eq:B}, \eqref{eq:C}, the
condition $N^*=\ln(P/d)<2$ (i.e. $d<P<de^2$) implies global attractivity of
$N^*$ for \emph{all} admissible delays and weights (Theorem~\ref{thm:main}),
matching the best delay-independent criterion for the autonomous equation
($N^*=2$ being the exact local-stability boundary); the monotone case
$N^*\le1$ is included.

\item[\textbf{(ii)}] The criterion $N^*<2$ survives intact when the constant
delay is replaced by an arbitrary variable delay $h(t)\in[t-s,t]$ (and by the
distributed/multiple/integral structures): Theorem~\ref{thm:main} needs no
condition on $h$. The delay-dependent (``$3/2$''-type) criteria valid for the
autonomous equation at $N^*\ge2$ for small $s$ do \emph{not} survive, since they
rest on delay smallness that large variable delays violate, and instability
occurs for $N^*>2$ once the delay exceeds $s_c$ (Corollary~\ref{cor:ii},
Proposition~\ref{prop:nec-auto}).

\item[\textbf{(iii)}] For the autonomous equation, uniform (in $s$) stability
requires $N^*\le2$; for $N^*>2$ the equilibrium is unstable for all $s>s_c$ with
$s_c$ in \eqref{eq:crit}. For each nonautonomous model, $N^*\le2$ is necessary
for stability uniform in the structure, and for a fixed structure realizing an
effective point delay $\tau$, stability requires \eqref{eq:necbound}: $\tau<s_c$
when $N^*>2$ (Proposition~\ref{prop:nec-nonauto}). Sufficient and necessary
delay-independent thresholds coincide at $N^*=2$, so the bound is sharp.
\end{itemize}

\end{document}
